% !TEX root = chapter-standalone.tex

\section{Functors generalize some familiar concepts}
\label{sec:functors-generalize}

In this section we collect further examples of functors.
We begin with four examples which show that functors generalize notions that we have already met: homomorphisms between graphs, \SY{monotone functions} between preorders, morphisms between monoids, and actions of monoids.

\subsection{Functors generalize graph homomorphisms}
\label{sec:functors-generalize-graph-homomorphisms}

Recall from \cref{sec:catsfromgraphs} that we can turn any graph $\graphA=\tupp{\vertices,\arcs, \source,\target}$ into a category $\pathcat{\graphA}$, the free category on $\graphA$ (\cref{def:free-category}).
The objects of $\pathcat{\graphA}$ are the vertices of the graph, and the morphisms from a vertex $\vertexa$ to a vertex $\vertexb$ are the paths from $\vertexa$ to $\vertexb$ (\cref{def:path}).
Composition is the concatenation of paths, and the identity morphisms are the trivial paths.

Now suppose that we have two graphs,
\begin{equation}
    {\graphA}_1=\tupp{\verticesn{1},\arcsn{1},\sourcen{1},\targetn{1}}
    \qqand
    {\graphA}_2=\tupp{\verticesn{2},\arcsn{2},\sourcen{2},\targetn{2}},
\end{equation}
and a \SY{graph homomorphism} $\funb \colon {\graphA}_1 \fto {\graphA}_2$ (\cref{def:graph-homomorphism}).
It consists of a function $\funbob \colon \verticesn{1} \sto \verticesn{2}$ on vertices and a function $\funbmor \colon \arcsn{1} \sto \arcsn{2}$ on edges, which are compatible with the source and target functions: an edge from $\vertexa$ to $\vertexb$ is mapped to an edge from $\funbob(\vertexa)$ to $\funbob(\vertexb)$.

From $\funb$ we can build a functor
\begin{equation}
    \funa \colon \pathcat{{\graphA}_1} \fto \pathcat{{\graphA}_2}
\end{equation}
as follows.
\begin{itemize}
    \item On objects, $\funa$ does what $\funb$ does on vertices: $\funaob \definedas \funbob$.
    \item On morphisms, $\funa$ applies $\funb$ to a path edge by edge:
          \begin{equation}
              \label{eq:graph-hom-functor-on-paths}
              \funamor\pars{\maketypedlist{e_1, \dots, e_n}{\arcsn{1}}} \definedas \maketypedlist{\funbmor(e_1), \dots, \funbmor(e_n)}{\arcsn{2}},
          \end{equation}
          and it maps the trivial path at a vertex $\vertexa$ to the trivial path at the vertex $\funbob(\vertexa)$.
\end{itemize}
The right-hand side of \cref{eq:graph-hom-functor-on-paths} is indeed a path in ${\graphA}_2$: since $\funb$ is compatible with sources and targets, the equation $\targetn{1}(e_i) = \sourcen{1}(e_{i+1})$ implies $\targetn{2}(\funbmor(e_i)) = \sourcen{2}(\funbmor(e_{i+1}))$.
For the same reason, a path from $\vertexa$ to $\vertexb$ is mapped to a path from $\funbob(\vertexa)$ to $\funbob(\vertexb)$, as the definition of a functor requires.
The two functor conditions hold as well: applying $\funb$ edge by edge to a concatenation of two paths gives the concatenation of the two resulting paths, and trivial paths are mapped to trivial paths by definition.

We have shown the following.

\begin{lemma}
    \label{lem:graph-hom-to-functor}
    Given graphs ${\graphA}_1$ and ${\graphA}_2$, every \SY{graph homomorphism} $\funb \colon {\graphA}_1 \fto {\graphA}_2$ induces a \SY{functor} $\funa \colon \pathcat{{\graphA}_1} \fto \pathcat{{\graphA}_2}$, which acts on vertices as $\funb$ does and which applies $\funb$ to paths edge by edge.
\end{lemma}

Different graph homomorphisms induce different functors, since we can recover $\funbob$ from $\funaob$, and $\funbmor$ from what $\funamor$ does on paths of length one.
So functors between free categories do generalize graph homomorphisms.
Note, however, that the correspondence is \emph{not} one-to-one: there are, in general, functors $\pathcat{{\graphA}_1} \fto \pathcat{{\graphA}_2}$ that are not induced by any graph homomorphism.

\begin{example}
    \label{exa:functor-not-from-graph-hom}
    Let $\graphA$ be the graph with a single vertex $\vertexa$ and a single edge $e$, which goes from $\vertexa$ to $\vertexa$.
    The morphisms of the category $\pathcat{\graphA}$ are the paths that go $n$ times around this loop, one for each $n \setin \natnumbers$, and composing two of them adds their lengths.

    There is a functor $\funa \colon \pathcat{\graphA} \fto \pathcat{\graphA}$ which is the identity on objects and which maps the path of length $n$ to the path of length $2n$.
    It is compatible with composition, since $2(n + m) = 2n + 2m$, and with identities, since $2 \cdot 0 = 0$.
    However, $\funa$ is not induced by any graph homomorphism $\graphA \fto \graphA$: a functor that is induced by a graph homomorphism maps paths of length one to paths of length one, whereas $\funa$ maps the path $\maketypedlist{e}{\arcs}$ to the path $\maketypedlist{e, e}{\arcs}$, which has length two.
\end{example}

In general, a functor between free categories is free to map each edge of the first graph to an arbitrary \emph{path} in the second graph, and the functors which are induced by graph homomorphisms are precisely those which map each path of length one to a path of length one.

In the next three subsections we will meet situations where the correspondence is even tighter: there, the functors between the categories in question correspond \emph{one-to-one} to the structure-preserving functions that we already know.

\todotext{J: the fact that not every functor between free categories is induced by a graph homomorphism is an interesting point to come back to when discussing adjunctions (free-forgetful adjunction between graphs and categories).}

\subsection{Functors generalize monotone functions}
\label{sec:posetsarecats}

\linkvideo{spring2021-tradeoffs:tradeoffs:orders:preorder-as-cat} % Pre-order as a category
\linkvideo{spring2021-tradeoffs:tradeoffs:orders:preorder-poset} % The skeleton of a pre-order is a poset
\linkvideo{spring2021-functors:semi-and-fun:mon-functions:mon-fun-as-func} % Monotone functions as functors

Recall from \cref{chap:structures-as-cats} that also a single \SY{preorder} $\posAdefinition$ can be viewed as a category $\poscat{\posA}$, called its categorification (\cref{def:poscat}).
The objects of $\poscat{\posA}$ are the elements of $\posAset$, and, given objects $\posela, \poselb \setin \posAset$, there is exactly one morphism $\posela \mto \poselb$ if $\posela \posAleq \poselb$, and there is no such morphism otherwise:
\begin{equation}
    \label{eq:poscat-homsets-recalled}
    \HomSet{\poscat{\posA}}{\posela}{\poselb} =
    \begin{cases}
        \singleton & \text{ if }  \ \posela \posAleq \poselb, \\
        \Emptyset  & \text{ else. }
    \end{cases}
\end{equation}
There is only one possible way to compose morphisms in $\poscat{\posA}$; that composites exist is guaranteed by the transitivity of $\posAleq$, and that identity morphisms exist is guaranteed by its reflexivity.

Now suppose that we have two preorders, $\posAdefinition$ and $\posBdefinition$.
What is a functor
\begin{equation}
    \funa \colon \poscat{\posA} \fto \poscat{\posB}
\end{equation}
between their categorifications?
Let us unpack \cref{def:functor} in this special case.
We will see that, in contrast to the situation for graphs, the answer is: exactly a monotone function.

The first constituent of $\funa$ is a function $\funaob \colon \Obof{\poscat{\posA}} \to \Obof{\poscat{\posB}}$.
Since the objects of $\poscat{\posA}$ and $\poscat{\posB}$ are the elements of $\posAset$ and $\posBset$, respectively, this is simply a function
\begin{equation}
    \funaob \colon \posAset \sto \posBset.
\end{equation}

The second constituent of $\funa$ is, for every pair of objects $\posela, \poselb \setin \posAset$, a function
\begin{equation}
    \label{eq:poscat-functor-on-morphisms}
    \funamor \colon \HomSet{\poscat{\posA}}{\posela}{\poselb} \to \HomSet{\poscat{\posB}}{\funaob(\posela)}{\funaob(\poselb)}.
\end{equation}
In view of \cref{eq:poscat-homsets-recalled}, there are two cases to consider.
\begin{itemize}
    \item If $\posela \posAleq \poselb$ does not hold, then the source of the function \cref{eq:poscat-functor-on-morphisms} is the empty set.
          There is exactly one function from the empty set to any given set, so in this case there is nothing to choose and nothing to check.
    \item If $\posela \posAleq \poselb$ holds, then the source of the function \cref{eq:poscat-functor-on-morphisms} is the set $\singleton$, which has one element.
          A function from $\singleton$ to the target set exists only if the target set is not empty, that is, only if
          \begin{equation}
              \funaob(\posela) \posBleq \funaob(\poselb).
          \end{equation}
          When this is the case, the target set is also $\singleton$, and there is exactly one function $\singleton \to \singleton$.
\end{itemize}
In summary: the functions $\funamor$ exist precisely when $\posela \posAleq \poselb$ implies $\funaob(\posela) \posBleq \funaob(\poselb)$ -- in other words, precisely when $\funaob$ is a \SY{monotone function} (\cref{def:monotone}) -- and, when they exist, there is no choice in how to define them.

Finally, the two conditions in \cref{def:functor} are equations between morphisms of $\poscat{\posB}$ that have the same source and the same target.
Since in $\poscat{\posB}$ there is at most one morphism with any given source and target, these equations hold automatically.

We have shown the following.

\begin{lemma}
    \label{lem:posetfunctor}
    Let $\posA$ and $\posB$ be preorders.
    \begin{enumerate}
        \item For every \SY{functor} $\funa \colon \poscat{\posA} \fto \poscat{\posB}$, the function $\funaob \colon \posAset \sto \posBset$ is a \SY{monotone function} $\posA \mto \posB$.
        \item For every \SY{monotone function} $\mapa \colon \posA \mto \posB$, there is exactly one \SY{functor} $\funa \colon \poscat{\posA} \fto \poscat{\posB}$ such that $\funaob = \mapa$.
    \end{enumerate}
    Hence the assignment $\funa \mapsto \funaob$ is a bijection between the set of functors $\poscat{\posA} \fto \poscat{\posB}$ and the set of monotone functions $\posA \mto \posB$.
\end{lemma}

It is in this sense that a monotone function $\posA \mto \posB$ is ``the same thing'' as a functor $\poscat{\posA} \fto \poscat{\posB}$.
We can read the lemma in two directions.
On the one hand, it says that functors are a generalization of monotone functions: monotone functions are what functors look like when the categories involved come from preorders.
On the other hand, it gives us an intuition for what the functor conditions mean in general.
A monotone function is one that ``respects the order''; a functor is one that ``respects composition''.
In a general category, however, there may be many different morphisms between two given objects, and so a functor must also specify \emph{which} morphism in the target category each morphism in the source category is sent to; and the two conditions in \cref{def:functor} are then no longer automatic.

\subsection{Functors generalize monoid morphisms}
\label{sec:functors-generalize-monoid-morphisms}

\linkvideo{spring2021-functors:semi-and-fun:ex-semigroup-semifun} % Semigroup morphisms as semi-functors

Recall from \cref{chap:structures-as-cats} that, similarly, a single \SY{monoid} $\monoidAdefinition$ can be viewed as a category $\monoidcat{\monoidA}$, called its categorification (\cref{def:monoidcat}).
This category has just one object, which we denote by $\monoidcatobj$.
Its morphisms are the elements of the monoid,
\begin{equation}
    \HomSet{\monoidcat{\monoidA}}{\monoidcatobj}{\monoidcatobj} = \monoidAset,
\end{equation}
composition of morphisms is given by the composition operation $\mtimesof\monoidA$ of the monoid, and the identity morphism of $\monoidcatobj$ is the neutral element $\idmon_\monoidA$.

Now suppose that we have two monoids, $\monoidAdefinition$ and $\monoidBdefinition$.
What is a functor
\begin{equation}
    \funa \colon \monoidcat{\monoidA} \fto \monoidcat{\monoidB}
\end{equation}
between their categorifications?
Again, let us unpack \cref{def:functor}.

The first constituent of $\funa$ is a function $\funaob \colon \Obof{\monoidcat{\monoidA}} \to \Obof{\monoidcat{\monoidB}}$.
Both categories have just the single object $\monoidcatobj$, so there is exactly one such function, and there is nothing to choose.

The second constituent of $\funa$ is a function
\begin{equation}
    \funamor \colon \HomSet{\monoidcat{\monoidA}}{\monoidcatobj}{\monoidcatobj} \to \HomSet{\monoidcat{\monoidB}}{\monoidcatobj}{\monoidcatobj},
\end{equation}
which is to say, a function
\begin{equation}
    \funamor \colon \monoidAset \sto \monoidBset.
\end{equation}

The first condition in \cref{def:functor}, compatibility with composition, says that for all $\monelA, \monelB \setin \monoidAset$,
\begin{equation}
    \label{eq:monoidcat-functor-composition}
    \funamor(\monelA \mtimesof\monoidA \monelB) = \funamor(\monelA) \mtimesof\monoidB \funamor(\monelB),
\end{equation}
and the second condition, compatibility with identity morphisms, says that
\begin{equation}
    \label{eq:monoidcat-functor-identity}
    \funamor(\idmon_\monoidA) = \idmon_\monoidB.
\end{equation}
These are precisely the two conditions \cref{eq:mon-mor-comp} and \cref{eq:mon-id-comp} that define a \SY{monoid morphism} (\cref{def:monoid-mor}).

We have shown the following.

\begin{lemma}
    \label{lem:monoid-hom-as-functor}
    Let $\monoidA$ and $\monoidB$ be monoids.
    \begin{enumerate}
        \item For every \SY{functor} $\funa \colon \monoidcat{\monoidA} \fto \monoidcat{\monoidB}$, the function $\funamor \colon \monoidAset \sto \monoidBset$ is a \SY{monoid morphism} $\monoidA \fto \monoidB$.
        \item For every \SY{monoid morphism} $\sgrpmorA \colon \monoidA \fto \monoidB$, there is exactly one \SY{functor} $\funa \colon \monoidcat{\monoidA} \fto \monoidcat{\monoidB}$ such that $\funamor = \sgrpmorA$.
    \end{enumerate}
    Hence the assignment $\funa \mapsto \funamor$ is a bijection between the set of functors $\monoidcat{\monoidA} \fto \monoidcat{\monoidB}$ and the set of monoid morphisms $\monoidA \fto \monoidB$.
\end{lemma}

It is in this sense that a monoid morphism $\monoidA \fto \monoidB$ is ``the same thing'' as a functor $\monoidcat{\monoidA} \fto \monoidcat{\monoidB}$.

\begin{example}
    In \cref{exa:string-length} we saw that the function which maps each string to its length is a \SY{monoid morphism} from a monoid of strings to the monoid $\tup{\natnumbers, +, 0}$.
    By \cref{lem:monoid-hom-as-functor}, we can equally well think of it as a functor between the corresponding one-object categories.
    Its compatibility with composition is the statement that the length of a concatenation of two strings is the sum of their lengths, and its compatibility with identities is the statement that the empty string has length $0$.
\end{example}

It is instructive to compare \cref{lem:posetfunctor} and \cref{lem:monoid-hom-as-functor}.
The categorification of a preorder may have many objects, but it has at most one morphism between any two of them.
Correspondingly, a functor between such categories is determined by what it does on objects, and the functor conditions hold automatically.
The categorification of a monoid, on the contrary, may have many morphisms, but it has only one object.
Correspondingly, a functor between such categories is determined by what it does on morphisms, and all of the content is in the functor conditions.
A functor between general categories combines both aspects.

\subsection{Functors generalize monoid actions}
\label{sec:functors-generalize-monoid-actions}

Let $\monoidAdefinition$ be a monoid.
Recall from \cref{def:monoid-cov-action} that an \SY{action} of $\monoidA$ on a set $\stateset$ is a \SY{monoid morphism}
\begin{equation}
    \act \colon \monoidA \fto \Endof\stateset,
\end{equation}
where $\Endof\stateset$ is the monoid whose elements are the functions $\stateset \to \stateset$, whose composition operation is the composition of functions, and whose neutral element is the identity function $\catidat\stateset$ (\cref{exa:endofunction-monoid}).

We will now see that monoid actions, too, are functors.
This time, we consider functors whose source is the categorification $\monoidcat{\monoidA}$ of the monoid, and whose target is the category $\Set$ of sets and functions.
So, what is a functor
\begin{equation}
    \funa \colon \monoidcat{\monoidA} \fto \Set?
\end{equation}
As before, we unpack \cref{def:functor}.

The first constituent of $\funa$ is a function $\funaob \colon \Obof{\monoidcat{\monoidA}} \to \Obof{\Set}$.
Since $\monoidcat{\monoidA}$ has just the single object $\monoidcatobj$, specifying such a function amounts to choosing one set, namely
\begin{equation}
    \stateset \definedas \funaob(\monoidcatobj).
\end{equation}

The second constituent of $\funa$ is a function
\begin{equation}
    \funamor \colon \HomSet{\monoidcat{\monoidA}}{\monoidcatobj}{\monoidcatobj} \to \HomSet{\Set}{\stateset}{\stateset}.
\end{equation}
The source of this function is the set $\monoidAset$, and its target is the set of functions $\stateset \to \stateset$.
In other words, $\funamor$ assigns to each element $\monelA \setin \monoidAset$ of the monoid a function
\begin{equation}
    \funamor(\monelA) \colon \stateset \sto \stateset.
\end{equation}

The first condition in \cref{def:functor} says that for all $\monelA, \monelB \setin \monoidAset$,
\begin{equation}
    \funamor(\monelA \mtimesof\monoidA \monelB) = \funamor(\monelA) \mthenof{\Set} \funamor(\monelB),
\end{equation}
and the second condition says that
\begin{equation}
    \funamor(\idmon_\monoidA) = \catidat\stateset.
\end{equation}
Since composition in $\Set$ is the composition of functions, these are precisely the conditions \cref{eq:act-to-endo-comp-compatibility} and \cref{eq:act-neutral-element-condit}, which say that $\funamor$ is a \SY{monoid morphism} from $\monoidA$ to $\Endof\stateset$.

We have shown the following.

\begin{lemma}
    \label{lem:monoid-action-as-functor}
    Let $\monoidA$ be a monoid.
    \begin{enumerate}
        \item For every \SY{functor} $\funa \colon \monoidcat{\monoidA} \fto \Set$, the function $\funamor$ is an action of $\monoidA$ on the set $\funaob(\monoidcatobj)$.
        \item For every set $\stateset$ and every action $\act \colon \monoidA \fto \Endof\stateset$ of $\monoidA$ on $\stateset$, there is exactly one \SY{functor} $\funa \colon \monoidcat{\monoidA} \fto \Set$ such that $\funaob(\monoidcatobj) = \stateset$ and $\funamor = \act$.
    \end{enumerate}
\end{lemma}

It is in this sense that a functor $\monoidcat{\monoidA} \fto \Set$ is ``the same thing'' as a set together with an action of $\monoidA$ on it.

\begin{example}
    Consider the flip-flop monoid $\mathbf{F}$ from \cref{exa:flip-flop-action}, and its action on the set $\makeset{0, 1}$ of possible contents of a memory cell.
    The corresponding functor $\funa \colon \monoidcat{\mathbf{F}} \fto \Set$ maps the single object $\monoidcatobj$ to the set $\makeset{0, 1}$.
    It maps the morphism $\text{set}$ to the function $\makeset{0, 1} \sto \makeset{0, 1}$ which is constant with value $1$, the morphism $\text{reset}$ to the function which is constant with value $0$, and the identity morphism $\idmon$ to the identity function.
\end{example}

\begin{remark}[Group actions]
    \label{rem:group-actions-as-functors}
    Every \SY{group} $\grpA = \tup{\grpAset, \mtimesof\grpA, \idgrp_\grpA, \ginv}$ is in particular a \SY{monoid}, so it, too, can be viewed as a category $\monoidcat{\grpA}$ with a single object $\monoidcatobj$, and \cref{lem:monoid-action-as-functor} applies: a functor $\funa \colon \monoidcat{\grpA} \fto \Set$ is the same thing as a set $\stateset$ together with a \SY{monoid morphism} $\funamor \colon \grpA \fto \Endof\stateset$.

    Recall, however, that an action of a group was defined as a group morphism $\grpA \fto \Autof\stateset$ to the group of invertible functions $\stateset \to \stateset$ (\cref{def:group-cov-action}).
    This is no discrepancy.
    For every $\grpelA \setin \grpAset$, the functor conditions and the inverse laws \cref{eq:group-inverse-law} give
    \begin{equation}
        \funamor(\grpelA) \mthenof{\Set} \funamor(\ginv(\grpelA)) = \funamor(\grpelA \mtimesof\grpA \ginv(\grpelA)) = \funamor(\idgrp_\grpA) = \catidat\stateset,
    \end{equation}
    and, in the same way, $\funamor(\ginv(\grpelA)) \mthenof{\Set} \funamor(\grpelA) = \catidat\stateset$.
    So each function $\funamor(\grpelA)$ is invertible, with inverse $\funamor(\ginv(\grpelA))$.
    In other words, $\funamor$ automatically takes its values in $\Autof\stateset$, and it is compatible with inverses.
    Hence a functor $\monoidcat{\grpA} \fto \Set$ is ``the same thing'' as a set together with an action of the group $\grpA$ on it.
\end{remark}

\paragraph{Set-valued functors as generalized actions.}
\Cref{lem:monoid-action-as-functor} suggests a useful point of view on functors whose target is the category $\Set$, whatever their source may be.
Let $\CatC$ be any category, and let
\begin{equation}
    \funa \colon \CatC \fto \Set
\end{equation}
be a functor.
Unpacking \cref{def:functor}, the functor $\funa$ provides
\begin{itemize}
    \item for every object $\Obja$ of $\CatC$, a set $\funaob(\Obja)$;
    \item for every morphism $\mora \colon \Obja \mtoin{\CatC} \Objb$, a function $\funamor(\mora) \colon \funaob(\Obja) \sto \funaob(\Objb)$.
\end{itemize}
Moreover, written out in terms of elements, the two functor conditions say that for all morphisms $\mora \colon \Obja \mtoin{\CatC} \Objb$ and $\morb \colon \Objb \mtoin{\CatC} \Objc$, and for all $\statesetel \setin \funaob(\Obja)$,
\begin{equation}
    \label{eq:generalized-action-conditions}
    \funamor(\mora \mthenof{\CatC} \morb)(\statesetel) = \funamor(\morb)\bigl(\funamor(\mora)(\statesetel)\bigr)
    \qqand
    \funamor(\catidat\Obja)(\statesetel) = \statesetel.
\end{equation}
These have exactly the same form as the conditions for a monoid action (compare \cref{eq:act-to-endo-comp-compatibility-elementwise}, and also \cref{def:uncurried-action}): acting by a composite $\mora \mthenof{\CatC} \morb$ is the same as first acting by $\mora$ and then acting by $\morb$, and acting by an identity morphism does nothing.

We can therefore think of any functor $\funa \colon \CatC \fto \Set$ as a \emph{generalized action}: an action of the category $\CatC$.
The difference to a monoid action is that the category $\CatC$ may have many objects.
Accordingly, there is not just one set that is acted upon, but a whole family of sets, one for each object of $\CatC$; and a morphism $\mora \colon \Obja \mtoin{\CatC} \Objb$ does not act on all of them, but only on the elements of $\funaob(\Obja)$, which it transforms into elements of $\funaob(\Objb)$.
If we think of the elements of the sets $\funaob(\Obja)$ as possible states of a system, then the objects of $\CatC$ play the role of different \emph{types} of states, and a morphism $\mora \colon \Obja \mtoin{\CatC} \Objb$ is a command which can only be applied to states of type $\Obja$, and which results in a state of type $\Objb$.
An action of a monoid is the special case where there is only one type.

We will meet several functors of this kind in what follows; see, for example, \cref{ex:GraphsViaFunctors}, where a very small category ``acts'' on the vertices and arrows of a graph.

\paragraph{Contravariant actions.}
What about \emph{contravariant} actions?
Recall that a contravariant action of a monoid $\monoidA$ on a set $\stateset$ is a \SY{monoid morphism} $\monoidA\op \fto \Endof\stateset$ whose source is the opposite monoid (\cref{def:monoid-cont-action}); in other words, it is an action of $\monoidA\op$.
Recall also that the categorification of the opposite monoid is the opposite of the categorification, $\monoidcat{\pars{\monoidA\op}} = \pars{\monoidcat{\monoidA}}\op$ (\cref{exa:opposite-cat-monoid}).
Applying \cref{lem:monoid-action-as-functor} to the monoid $\monoidA\op$, we find that a contravariant action of $\monoidA$ is ``the same thing'' as a functor
\begin{equation}
    \pars{\monoidcat{\monoidA}}\op \fto \Set.
\end{equation}

Accordingly, for any category $\CatC$, we can think of a functor $\funa \colon \CatCop \fto \Set$ as a generalized contravariant action of $\CatC$.
Unpacking the definition of the opposite category (\cref{def:opposite-category}), such a functor provides a set $\funaob(\Obja)$ for every object $\Obja$ of $\CatC$, and, for every morphism $\mora \colon \Obja \mtoin{\CatC} \Objb$, a function in the reverse direction,
\begin{equation}
    \funamor(\mora\op) \colon \funaob(\Objb) \sto \funaob(\Obja).
\end{equation}
Acting by a composite $\mora \mthenof{\CatC} \morb$ is now the same as first acting by $\morb$ and then acting by $\mora$ (compare \cref{eq:uncurried-cont-action-1cond}).
An example which combines both kinds of action is the hom-functor of \cref{exa:hom-functor}.


\paragraph{Canonical actions of a category.}


\begin{example}\label{exa:hom-functor}
    Let $\CatC$ be any category.
    There is a functor
    \begin{equation}\label{eq:hom-functor-type}
        \funa \colon \CatC\op \Ctimes \CatC \fto \Set
    \end{equation}
    defined on objects by
    \begin{equation}\label{eq:hom-functor-on-objects}
        \funaob(\tup{\Obja, \Objb}) = \Hom_\CatC(\Obja, \Objb).
    \end{equation}
    To see how $\funa$ is defined on morphisms, let $\mora\op \colon \Objd \mtoin{\CatC\op} \Obja$ and $\morb \colon \Objb \mtoin{\CatC} \Objc$ be morphisms in $\CatC\op$ and $\CatC$, respectively, and recall that $\mora\op \colon \Objd \mtoin{\CatC\op} \Obja$ is, by definition, a morphism $\mora \colon \Obja \mtoin{\CatC} \Objd$.
    We define $\funamor(\tup{\mora\op, \morb})$ to be the function
    \begin{equation}\label{eq:hom-functor-on-morphisms}
        \defmapperiod{\funamor(\tup{\mora\op, \morb})}{\Hom_\CatC(\Obja, \Objb)}{\sto}{\Hom_\CatC(\Objd, \Objc)}{\stylemorph{\varphi}}{\mora \mthen \stylemorph{\varphi} \mthen \morb}
    \end{equation}



